% Compile with: uplatex file.tex && dvipdfmx file.dvi
%           or: latexmk -uplatex file.tex
\documentclass[dvipdfmx]{article}
\usepackage{amsmath, amssymb}
\usepackage{tikz}

\title{scenario\_central\_force\_conservation}
\date{\today}

\begin{document}
\maketitle

Consider a point mass $m$ moving in a central-force potential in polar coordinates $(r,\theta)$. $k$ is a parameter giving the strength of the Kepler potential $U=-k/r$.

Define r as a time-dependent function

$\checkmark$ Declared r as time-dependent (time variable: t)

Define theta as a time-dependent function

$\checkmark$ Declared theta as time-dependent (time variable: t)

Define m, k as scalars

\[
  m, k
\]

Define the angular momentum $L_z=mr^2\dot\theta$ (mass times the square of the radius times the angular velocity).

$L_z = m r^2 \dot{\theta}$

\begin{equation}
  L_z = m r^2 \dot{\theta}
  \label{eq:E1}
\end{equation}

To check whether the angular momentum is conserved, take the total time derivative of $L_z$ (the total derivative includes the chain rule, accounting for the time dependence of both $r$ and $\theta$).

Take the total derivative of E1 with respect to t

\begin{equation}
  \frac{dL_z}{dt} = m r^2 \left(\frac{d^2}{d t^2} \theta\right)+2 m r \dot{r} \left(\frac{d}{d t} \theta\right)
  \label{eq:E2}
\end{equation}

Substitute the equation of motion $\ddot\theta=-2\dot r\dot\theta/r$ into $dL_z/dt$. If the result is zero, the angular momentum is a conserved quantity.

In E2, substitute -2 $\dot{r}$ $\dot{\theta} / r$ for $\ddot{\theta}$

\begin{equation}
  \frac{dL_z}{dt} = 0
  \label{eq:E3}
\end{equation}

Define the energy $E=\frac12 m(\dot r^2+r^2\dot\theta^2)-k/r$ (the sum of the kinetic energy and the Kepler potential).

$E_tot = (1/2) m (\dot{r}^2 + r^2 \dot{\theta}^2) - k / r$

\begin{equation}
  E_tot = (1/2) m (\dot{r}^2 + r^2 \dot{\theta}^2) - k / r
  \label{eq:E4}
\end{equation}

To check whether the energy is conserved, take the total time derivative of $E$.

Take the total derivative of E4 with respect to t

\begin{equation}
  \frac{dE_tot}{dt} = \frac{m \left(2 r^2 \left(\frac{d}{d t} \theta\right) \left(\frac{d^2}{d t^2} \theta\right)+2 r \dot{r} \left(\frac{d}{d t} \theta\right)^2+2 \dot{r} \ddot{r}\right)}{2}+\frac{k \dot{r}}{r^2}
  \label{eq:E5}
\end{equation}

Substitute the equation of motion $\ddot\theta=-2\dot r\dot\theta/r$.

In E5, substitute -2 $\dot{r}$ $\dot{\theta} / r$ for $\ddot{\theta}$

\begin{equation}
  \frac{dE_tot}{dt} = -\frac{m r^3 \dot{r} \left(\frac{d}{d t} \theta\right)^2-m r^2 \dot{r} \ddot{r}-k \dot{r}}{r^2}
  \label{eq:E6}
\end{equation}

Substitute the other equation of motion $\ddot r=r\dot\theta^2-k/(mr^2)$ as well. If the result of substituting both is zero, the energy is a conserved quantity.

In E6, substitute r $\dot{\theta}^2 - k$ / (m $r^2$) for $\ddot{r}$

\begin{equation}
  \frac{dE_tot}{dt} = 0
  \label{eq:E7}
\end{equation}

\end{document}